%------------------------
% Resume Template
% Author : Anubhav Singh
% Github : https://github.com/xprilion
% License : MIT
%------------------------

\documentclass[a4paper]{article}

\usepackage{latexsym}
\usepackage[empty]{fullpage}
\usepackage{titlesec}
\usepackage{marvosym}
\usepackage[usenames,dvipsnames]{color}
\usepackage{verbatim}
\usepackage{enumitem}
% \usepackage[pdftex]{hyperref}
\usepackage{fancyhdr}

\usepackage[hidelinks]{hyperref}

\usepackage{fontawesome5}
\makeatletter
\newcommand{\github}[1]{%
   \href{#1}{\faGithubSquare}%
}
\makeatother
\pagestyle{fancy}
\fancyhf{} % clear all header and footer fields
\fancyfoot{}
\renewcommand{\headrulewidth}{0pt}
\renewcommand{\footrulewidth}{0pt}

% Adjust margins
\addtolength{\oddsidemargin}{-0.530in}
\addtolength{\evensidemargin}{-0.375in}
\addtolength{\textwidth}{1in}
\addtolength{\topmargin}{-.60in} % original
% \addtolength{\topmargin}{-.80in} % 80 for sci and engg library
\addtolength{\textheight}{1in}

\urlstyle{rm}

\raggedbottom
\raggedright
\setlength{\tabcolsep}{0in}

% Sections formatting
\titleformat{\section}{
  \vspace{-10pt}\scshape\raggedright\large
}{}{0em}{}[\color{black}\titlerule \vspace{-6pt}]

%-------------------------
% Custom commands
% \newsavebox\CBox
% \def\textBF#1{\sbox\CBox{#1}\resizebox{\wd\CBox}{\ht\CBox}{\textbf{#1}}}
% \parindent=0pt

% \newcommand{\resumeItem}[2]{
%   \item\small{
%     \textBF{#1}{: #2 \vspace{-2pt}}
%   }
% }
% \newcommand{\resumeItem}[2]{
%   \item\small{
%     {#1}{: #2 \vspace{-2pt}}
%   }
% }
\newcommand{\resumeItem}[2]{
  \item\small{
    \textbf{#1}{: #2 \vspace{-2pt}}
  }
}

\newcommand{\resumeItemWithoutTitle}[1]{
  \item\small{
    {\vspace{-2pt}}
  }
}

\newcommand{\resumeSubheading}[4]{
  \vspace{-1pt}\item
    \begin{tabular*}{0.97\textwidth}{l@{\extracolsep{\fill}}r}
      \textbf{#1} & #2 \\
      {#3} & {#4} \\
    \end{tabular*}\vspace{-5pt}
}

\newcommand{\resumeSubheadingSecondHalf}[2]{
  \vspace{-1pt}
    \begin{tabular*}{0.97\textwidth}{l@{\extracolsep{\fill}}r}
      % \textbf{#1} & #2 \\
      {#1} & {#2} \\
    \end{tabular*}\vspace{-5pt}
}


\newcommand{\resumeSubItem}[2]{\resumeItem{#1}{#2}\vspace{-3pt}}

\renewcommand{\labelitemii}{$\circ$}

\newcommand{\resumeSubHeadingListStart}{\begin{itemize}[leftmargin=*]}
\newcommand{\resumeSubHeadingListEnd}{\end{itemize}}
\newcommand{\resumeItemListStart}{\begin{itemize}[leftmargin=*]}
% \newcommand{\resumeItemListStart}{\begin{itemize}[leftmargin=0.2cm]}
% \newcommand{\resumeItemListStart}{\begin{itemize}}
\newcommand{\resumeItemListEnd}{\end{itemize}\vspace{-5pt}}

%-----------------------------
%%%%%%  CV STARTS HERE  %%%%%%

\begin{document}

%----------HEADING-----------------
\centering{\textbf{\underline{\LARGE Ojas Jeetendra Raundale}}} \\
\begin{tabular*}{\textwidth}{l@{\extracolsep{\fill}}r}
  \href{https://www.linkedin.com/in/ojasraundale/}{linkedin.com/in/ojasraundale/} & (413)-328-5823 \\
  \href{https://github.com/ojasraundale}{github.com/ojasraundale} & \href{mailto:}{raundale.ojas@gmail.com}\\
\end{tabular*}

%-----------EDUCATION-----------------
\section{Education}
  \resumeSubHeadingListStart
    \resumeSubheading
      {University of Massachusetts Amherst}{Massachusetts, US}
      {Masters of Science (MS) - Computer Science; GPA: 3.87/4.0  }{Sept 2023 - May 2025 (exp.)}
      % \vspace{1pt}
      {\scriptsize { \footnotesize{\newline{}\textbf{Courses:} Reinforcement Learning,
      Natural Language Processing,
      Probabilistic Graphical Models,
      % Systems for Data Science,
      Big Data,
      Distributed Systems
      % Advanced Natural Language Processing - Large Language Models,
      % Deep Learning for 3D Vision,
      % Quantum Information Systems 
      }}}
      \vspace{-4pt}
    \resumeSubheading
      {Indian Institute of Technology Dharwad}{Karnataka, India}
      {Bachelor of Technology - Computer Science and Engineering; 
      % GPA: 8.68/10
      }{July 2017 - June 2021}
      % {\scriptsize { \footnotesize{\newline{}\textbf{Relevant Courses:} Machine Learning \& Neural Networks, Convex Optimization Theory, Data Structure \& Algorithms, Data Analysis}}}

% \vspace{-10pt}

      
% Machine Learning \& Pattern Recogn.\\
% Convex Optimization Theory \\
% Data Structure and Algorithms\\
% Probability and Random Processes\\
% Graph Theory \& Combinatorics\\
% Database and Information Systems\\
% Computer Architecture\\
% Operating Systems\\
% Compilers\\
% Computer Networks\\
% Software Systems\\
% Data Analysis\\
% Discrete Structures \& Automata Theory\\
% Design and Analysis of Algorithms\\
% Parameterized Algorithms\\

    \resumeSubHeadingListEnd
	    
\vspace{1pt}

\section{Experience}
  \resumeSubHeadingListStart


\resumeSubheading{Fidelity Investments}{Boston, MA}
    {\underline{Data Scientist}}{Jun 2025 - Present}
    \resumeItemListStart
        \resumeItemWithoutTitle{Edubot (Chat with my Data) / fine-tuning}
        % {Developed a text-to-structured-query pipeline by \textbf{fine-tuning LLMs} using SFT/RLHF for \textbf{JQ (JSON Query) prediction}, streamlining natural language access to complex hierarchical datasets.}
        % Developed a structured query generation pipeline for the "Chat with my Data" flow by fine-tuning LLMs (SFT, RLHF) to generate schema-aware JQ filters over complex hierarchical data, enabling natural language querying of structured datasets.
        Developed a natural-language-to-structured-query generation pipeline for the "Chat with my Data" flow by \textbf{fine-tuning LLMs (SFT, RLHF)} to generate schema-aware JQ filters, enabling targeted extraction of relevant subsets from large customer datasets.
        
        
        \resumeItemWithoutTitle{Repeatability}
        {
        % Quantified model non-determinism by implementing \textbf{distribution-based population metrics} and entropy analysis to evaluate response consistency, significantly reducing variance in JSON Query generation.
        % Developed a repeatability evaluation framework by performing multi-sample inference per input and analyzing output distributions using entropy and population-level metrics; leveraged these signals to quantify non-determinism and improve the LLM systems.
        Designed a \textbf{repeatability-focused evaluation} framework using entropy-based and population-level distributional metrics to \textbf{quantify non-determinism} and drive improvements in output consistency for LLM-based systems.
        }
        
        
        % \resumeItemWithoutTitle{Advanced Evaluation Frameworks}
        % {Built a multi-layered benchmarking suite using \textbf{LLM-as-a-Judge} and character-level \textbf{Span metrics}, providing granular accuracy tracking for intent classification and entity highlighting modules.}


        \resumeItemWithoutTitle{Agentic AI and Tool Calling}
        {
        % Improved LLM tool-calling for agentic workflows by building benchmarks, synthetic datasets, and evaluation metrics; compared models across accuracy, latency, and throughput for selection, and refined tool descriptions/prompts to ensure consistent behavior in production.
        Improved LLM \textbf{tool-calling} for \textbf{agentic workflows} by designing evaluation benchmarks and grounded \textbf{synthetic datasets} to compare generated tool calls against gold standards; evaluated models across accuracy, latency, and throughput, and refined tool definitions and prompts for reliable production behavior.
        }
        
        % \resumeItemWithoutTitle{Data Engineering & Annotation}
        % {Spearheaded the creation of "Gold" training sets by orchestrating large-scale human-in-the-loop workflows, utilizing \textbf{Inter-Annotator Agreement (IAA)} to resolve labeling conflicts and ensure high-fidelity ground truth.}
    \resumeItemListEnd

%     \resumeSubheading{Fidelity Investments}{Boston, MA}
% {\underline{Data Scientist}}{Jun 2025 - Present}
% \resumeItemListStart

% \resumeItem
% {Built and improved LLM-based systems for natural language interaction with structured data, including a text-to-query pipeline using SFT and RLHF for JSON query generation.}

% \resumeItem
% {Improved LLM tool-calling reliability for agentic workflows by building custom benchmarks, synthetic datasets, and evaluation metrics; compared models across accuracy, latency, and throughput for selection, and refined tool descriptions/prompts to ensure consistent behavior in production.}

% \resumeItem
% {Developed evaluation and reliability frameworks using LLM-as-a-Judge, span-level metrics, and entropy-based analysis to quantify non-determinism and track fine-grained model performance across tasks.}

% \resumeItemListEnd
        
  
\resumeSubheading{Google DeepMind | IESL Lab, University of Massachusetts, Amherst}{Remote, USA}
    {\underline{Graduate Research Extern}}{Jan 2025 - May 2025}
    \resumeItemListStart
\vspace{1pt}
        \resumeItemWithoutTitle{Multi Touch Attribution Project}
    {
         Researched and developed \textbf{Quadratic Model Ensembling} methods for LLMs and VLMs using adaptive optimization techniques in collaboration with \textbf{Google DeepMind}
    }
        \resumeItemWithoutTitle{Multi Touch Attribution Project}
    {
        Investigating Out-of-Distribution performance, and computational expenditure of Ensembling Methods in comparison with fine-tuning and Model SOUPing. 
    }
        \resumeItemListEnd

          
  \resumeSubheading{IBM | IESL Lab, University of Massachusetts, Amherst}{Amherst, MA}
    {\underline{Graduate Research Extern}}{August 2024 - May 2025}
    \resumeItemListStart
\vspace{1pt}
      \resumeItemWithoutTitle{Multi Touch Attribution Project}
    {
         Developing and researching \textbf{fine-tuning algorithms for LLMs} based on distribution matching methods under Prof. Andrew McCallum and in collaboration with \textbf{IBM Research}.
    }
      \resumeItemWithoutTitle{Multi Touch Attribution Project}
    {
        Utilised concepts from \textbf{Markov Chains} and \textbf{Self-Play Optimizations} to achieve accuracies better than Direct Preference Optimization (DPO) and Reinforcement learning from human feedback (RLHF).
    }
          \resumeItemListEnd
  
  \resumeSubheading{FCAT, Fidelity Investments}{Boston, MA}
    {\underline{Data Science Intern}}{June 2024 - August 2024}
    \resumeItemListStart
\vspace{1pt}
      \resumeItemWithoutTitle{Multi Touch Attribution Project}
    {
        Reduced ChatGPT API costs by more than 95\% by replacing it with in-house \textbf{fine-tuned Large Language Models}. 
        % Crafted a POC to save more than 95\% on API costs by replacing ChatGPT with in-house, open-sourced \textbf{fine-tuned} \textbf{LLMs}. 
    }
      \resumeItemWithoutTitle{Multi Touch Attribution Project}
    {
        Automated the process of support-customer phone call evaluation by fine-tuning  \textbf{Phi-3} and \textbf{Llama 3.1} using \textbf{LORA} and \textbf{qLORA} and achieving accuracies greater than 90\%. 
        % Minimized the manual component of quality control in support-customer phone calls by fine-tuning \textbf{Phi-3} and \textbf{Llama 3.1} models to autonomously evaluate and rate support agents' performance.
    }
                
      \resumeItemWithoutTitle{Multi Touch Attribution Project}
    {
          % Deployed a \textbf{cost-effective, scalable architecture} to host the LLM system on AWS Lambda Clusters.
          Designed custom LLM evaluation metrics and utilized a multi-GPU distributed system to train LLMs using \textbf{Deepspeed}. %and \textbf{Accelerate}.
          % Achieved accuracies greater than 90\% by using \textbf{LORA} and \textbf{qLORA}, utilizing decoding techniques, designing custom LLM evaluation metrics, and training on a multi-GPU distributed system using \textbf{Deepspeed} and \textbf{Accelerate}.
      }          
          \resumeItemListEnd
          
    \resumeSubheading{Express Analytics}{Pune, India}
    {\underline{Machine Learning Engineer - Data Science}}{June 2021 - June 2023}
    \resumeItemListStart
\vspace{1pt}
      %   \resumeItem{Multi Touch Attribution}
      %     {Owner of the Multi-Touch Attribution Module: Designed Deep Learning and Markov based MTA Models in Python's Keras. The model looks at customers' journeys from Google Analytics Data to figure out the optimal marketing spend that the client needs to assign to different channels and campaigns.}
      %     \resumeItem{Customer Life Time Value Prediction}
      %     {Developed Deep Learning-based solutions to predict the lifetime value of customers. Enhanced the segmenting of customers to plan out the target audiences for marketing campaigns.}
      %     \resumeItem{Continuous Integration \& Deployment}{Revamped the deployment procedures by creating CI/CD pipelines and Dockerized microservices on Gitlab and Azure DevOps reducing development time and effort.}
      %     \resumeItem{Visualization and Analytics}{Developed and maintained dashboards for clients using Power Bi, Tableau and Apache Superset increasing interactivity between stakeholders and developers.}
      %     \resumeItem{A/B testing}{Experimented over 100 hypotheses monthly using R to find marketing decisions with the most impact. }
      % \resumeItemListEnd
      \resumeItemWithoutTitle{Multi Touch Attribution Project}
        {Improved customer targeting and marketing effectiveness resulting in increased ROI by developing custom \textbf{Multi-Touch Attribution and Marketing Mix Models.}}
      \resumeItemWithoutTitle{Multi Touch Attribution Project}
      {
      Analyzed millions of customer journeys from \textbf{Google Analytics Data} by employing models built upon RNNs, LSTMs, and Statistical Models in \textbf{Python's Keras} to optimize marketing campaign spends.
      }
      % {Owned the development of two major projects, employing \textbf{RNNs}, \textbf{LSTMs}, and \textbf{Statistical models} in Python’s \textbf{Keras} to optimize marketing spend for clients.}
          
      % \resumeItemWithoutTitle{Customer Life Time Value Prediction}
          % {Improved \textbf{customer targeting and segmentation} for marketing campaigns using Deep Learning-based solutions to predict \textbf{customer lifetime value}, leading to more effective marketing campaigns.}
      % {Analyzed millions of customer journeys from \textbf{Google Analytics Data} and implemented custom Multi-Touch Attribution models to determine the most effective budget allocation, resulting in improved marketing, customer targeting and ROI.}
    
      \resumeItemWithoutTitle{Continuous Integration \& Deployment}{Revamped the deployment procedures by creating \textbf{CI/CD pipelines} and \textbf{Dockerized microservices} on \textbf{Jenkins} DevOps and \textbf{Gitlab} reducing deployment effort from days to hours.}
      
      % \resumeItemWithoutTitle{A/B testing}
      % {Conducted 100+ monthly \textbf{A/B} tests using \textbf{R} to evaluate marketing decisions providing valuable insights and contributing to data-backed adjustments.}
          % \resumeItemWithoutTitle{Dashboard Development and Training Initiative}{Designed \textbf{dashboards} for clients using \textbf{Power Bi, Tableau and Apache Superset}. Also conducted a company-wide workshop for developers to get started with using PowerBi. Boosted client satisfaction and promoted a data-driven, self-service analytics culture within the organization.}
          % \resumeItemWithoutTitle{Dashboard Development and Training Initiative}{Designed \textbf{dashboards} for clients using \textbf{Power Bi, and Tableau}. Also conducted a company-wide workshop for developers to get started with using PowerBi. Boosted client satisfaction and promoted a data-driven, self-service analytics culture within the organization.}
          % \resumeItemWithoutTitle{Dashboard Development and Training Initiative}{Designed client-side \textbf{dashboards} using \textbf{Power BI} and \textbf{Tableau}. Conducted a company-wide Power BI workshop, boosting client satisfaction and promoting a data-driven, self-service analytics culture.}
      % \resumeItemWithoutTitle{Dashboard Development and Training Initiative}{Designed client dashboards using \textbf{Power BI} and \textbf{Tableau} to boost client satisfaction and conducted a company-wide Power BI workshop to promote a data-driven, self-service analytics culture.}
          
      \resumeItemListEnd
    % \vspace{-2pt}
%       \resumeSubheadingSecondHalf
%     {\underline{Data Science Intern}}{May 2020 -  Aug 2020}
%     % \vspace{1pt}
%     \resumeItemListStart
% \vspace{1pt}
%         \resumeItemWithoutTitle{Markov Based MTA modeling}
%           {Engineered a data-dependent \textbf{multi-touch attribution} model utilizing \textbf{Markov} based approaches to address limitations of existing single-touch attribution methods and providing attributions at a customer level.  }
%     \resumeItemListEnd

    % \resumeItemListEnd
  %   \resumeSubheading
		% {Express Analytics}{Internship}
		% {Data Science Intern}{May 2020 -  Aug 2020}
		% \resumeItemListStart
  %       \resumeItem{Multi Touch Attribution}
  %         {Attribution modeling for a client's website data to assign different budgets for the different marketing channels which the users use to arrive at their website. Designed a Markov Model-based attribution model which upon looking at the users' activity and whether the user purchased a product or not, assigns attributions to the marketing channels which the user took.}
  %         \resumeItemListEnd
% \vspace{-3pt}
  %   \resumeSubheading
		% {MSCI Inc.}{Internship}
		% {\underline{Java Intern}}{Feb 2021 -  May 2021}
		% % \resumeItemListStart
  % %       \resumeItem{Java module to run python scripts on the backend}
  % %         {Accelerated the research team's efforts to find the impact of python based POC functions before they were implemented in Java backend.}
		% % \resumeItemListEnd
  %           \resumeItemListStart
  %       \resumeItemWithoutTitle{}
  %         {Developed a \textbf{Java} module that enabled the execution of Python scripts on the backend, expediting the research team's evaluation of Python-based POC functions within the Java system. This streamlined process enhanced decision-making and accelerated software development efforts.}
		% \resumeItemListEnd
% \vspace{-3pt}
    


  
        %   \resumeItem{}
        % %   {Designed a Markov Model based attribution model which upon looking at the users activity and whether the user purchased a product or not, assigns attributions to the marketing channels which the user took.}
       
% \vspace{-3pt}
%     \resumeSubheading
% 		{Indian Institute of Technology, Bombay}{Internship}
% 		{\underline{NLP Intern}}{May 2019 -  July 2019}
% 		\resumeItemListStart
% \vspace{1pt}
%         \resumeItemWithoutTitle{Entity Relationship Extraction using NLTK in python}
%           % {Designed a \textbf{Python} module using NLTK toolkit and Stanford CoreNLP API that produces Entity Relationships for a given paragraph describing a database. The objective was to create a website just out of a descriptive paragraph. The ER Model was an essential component used to create the basic meta-model of the database upon which the website could be built.}
%           % {Designed a \textbf{Python} module using NLTK toolkit and Stanford CoreNLP API to produce Entity Relationships for a given paragraph describing a database. The module generates The objective was to create a website just out of a descriptive paragraph. The ER Model was an essential component used to create the basic meta-model of the database upon which the website could be built.}
%           {
%           Contributed to the goal of creating a website from only a descriptive paragraph by designing a \textbf{Python} module that creates a meta-model of the underlying database. 
%           }
%           \resumeItemWithoutTitle{Entity Relationship Extraction using NLTK in python}
%           {
%             Utilized \textbf{NLTK} toolkit, and \textbf{Stanford CoreNLP} to extract Entity Relation features - an essential part of the meta-model. 
%           }
		% \resumeItemListEnd
\resumeSubHeadingListEnd

%-----------PROJECTS-----------------
\vspace{-8pt}
\section{Projects}
\resumeSubHeadingListStart
% \resumeSubheadingSecondHalf{Crop Classification using Multispectral drone images (Image processing, CNN, Deep Learning)}{2020}

\resumeSubItem{Medical Minds Unite - Under guidance of Prof. Mohit Iyyer (NLP, Transformers, PyTorch)
\href{https://github.com/ojasraundale/medical-minds-llm}{\faGithubSquare}
}
{
\newline Strengthened the performance of open-source LLMs for medical applications by fine-tuning, quantizing, and implementing merging techniques such as Ensembling, Mixture of Experts, and QLora.
} 

\resumeSubItem{Statistical Models Repository (JAX, PyTorch, Monte Carlo Methods)}
{\href{https://github.com/stars/ojasraundale/lists/statistical-models}{\faGithubSquare}
\newline 
% A growing list of statistical model implementations to explore off-trend methods that solve complex problems such as Generation, and Bradley-Terry Ranking. Experimented with Hamiltonian Markov Chain Monte Carlo Methods, Gibbs Sampling, Message Passing, Conditional Random Fields, and Variational Autoecoders to achieve a satisfactory performance. 
Developed a repository of statistical model implementations to explore innovative methods for solving complex problems such as Generation and Bradley-Terry Ranking. Experimented with Hamiltonian MCMC, Message Passing, Conditional Random Fields, and Variational Autoencoders resulting in robust and effective solutions.
% Some key projects - 1.) Hamiltonian Markov Chain Monte using Gibbs Sampling in Python's JAX to model the Bradley-Terry ranking problem. 2.) An implementation of Message Passing and Learning of Conditional Random Fields to learn english words in an image dataset.  
} 
\resumeSubItem{Reinforcement Learning Repository (Deep Learning, Reinforcement Learning, PyGame, PyTorch)}{\href{https://github.com/stars/ojasraundale/lists/reinforcement-learning}{\faGithubSquare}
\newline Assembled a Portfolio of Reinforcement Learning Algorithm implementations like Policy Gradient methods, Q-learning, SARSA, and Evolution strategies. Trained and evaluated RL agents on classical RL domains such as cart pole and lunar lander as well as self-made games helping the AI bot play the game with optimal performance.}


\vspace{2pt}
% \resumeSubItem{Crop Classification using Multispectral drone images (Image processing, CNN, Deep Learning)}{\href{https://github.com/ojasraundale/drone-crop-image-analysis}{\faGithubSquare}
% \newline Constructed a Deep Learning model to predict the crop cover, plant health, and species of different vegetation present in multi-spectral drone images containing 5 different bands of light. Applied photogrammetry techniques using Micasense to create orthomosaics of the field.}



% % \vspace{2pt}
% \resumeSubItem{Pet Activity Change Detection (Python, Hidden Markov Models, Deep Learning)}{
% \href{https://github.com/ojasraundale/pet-activity-change-detection}{\faGithubSquare}
% \newline Implemented K-means, Gaussian Hidden Markov Model, and LSTM Neural Networks to detect change in the activities of pet animals based on timed accelerometer and gyroscopic data read from sensor belts.}


% % \vspace{2pt}
% \resumeSubItem{MICRO Compiler (C, Flex, Bison, Git)}{\href{https://github.com/ojasraundale/micro-compiler}{\faGithubSquare}
% \newline Programmed a compiler for MICRO language that included scanning, parsing, symbol table formation, semantic actions, optimizations and register allocations using C-based Flex and Bison utilities.}
% % \vspace{2pt}



% \vspace{2pt}
% \resumeSubItem{Simulation of a Processor (Java) }
% {\newline
% Simulated a pipelined processor in Java. Given an input program in assembly language, the processor executes it's code and produces output. Analyzed the change in performance of the processor on a set of benchmark programs with different types of caches and pipelines.}

\resumeSubHeadingListEnd
% \vspace{-10pt}



% %-----------PROJECTS-----------------
% \vspace{-5pt}
% \section{Projects}
% \resumeSubHeadingListStart

% \resumeSubheading{Crop Classification using Multispectral drone images}{Final year B.Tech Project}
%     {Image processing, CNN, Deep Learning}{Aug - Dec 2020}
%     { {\small{\newline{} Developed a DL-based model to predict the crop cover, plant health and species of different vegetation present in multi-spectral drone images. The images contained 5 different bands of light (RGB + Near Infra-Red and Red-edge). Various photogrammetry techniques using Micasense and Agisoft Metashape to create orthomosaics of the field.}}}
    
% \vspace{0pt}
% \resumeSubheading{MICRO Compiler}{Compilers Course Project}
%     {C, Flex, Bison, Git}{Jan - Apr 2020}
%      {\small{\newline{} {Implemented a compiler for MICRO language. The build of the compiler includes scanning, parsing, symbol table formation, semantic actions, optimizations, and register allocations. It used C-based Flex and Bison utilities}}}
    
% \vspace{0pt}
% \resumeSubheading{Pet Activity Change Detection }{PRML Course Project}
%     {Python, Hidden Markov Models, Deep Learning}{Aug - Nov 2019}
%      {\small{\newline{} {Implemented K-means, Gaussian Hidden Markov Model, and LSTM Neural Networks to detect change in the activities of pet animals based on accelerometer and gyroscopic data. The data came from sensor belts attached to pets and recorded at different instances of time.}}}

% \vspace{0pt}
% \resumeSubheading{Simulation of a Processor}{Computer Architecture Course Project}
%     {Java, OOPS}{Aug - Nov 2019}
%     { {\small{\newline{} {Simulated a pipelined processor in Java. Given an input program in assembly language, the processor will execute the code and produce the results. Analyzed the change in performance of the processor on a set of benchmark programs with different types of caches and pipelines.}}}}

% \vspace{0pt}
% \resumeSubheading{Bulls And Cows Number Guesser}{Data Structures Course Project}
%     {C++, Bash scripting, OOPS}{Jul - Nov 2018}
%     { {\small{\newline{} {Developed a C++ program to guess the "Secret" Number of the opponent team in the game of Bulls \& Cows (A variation of Mastermind). Implemented Min-Max Algorithm and achieved an average of 5.24 moves of guessing the secret number. }}}}
% \resumeSubHeadingListEnd
% \vspace{-5pt}

%  % %-----------PROJECTS-----------------
% \vspace{-5pt}
% \section{Projects}
% \resumeSubHeadingListStart

% \resumeSubheading{Crop Classification using Multispectral drone images}{Final year B.Tech Project}
%     {Image processing, CNN, Deep Learning}{Aug - Dec 2020}
%     { {\small{\newline{} Developed a DL-based model to predict the crop cover, plant health and species of different vegetation present in multi-spectral drone images. The images contained 5 different bands of light (RGB + Near Infra-Red and Red-edge). Various photogrammetry techniques using Micasense and Agisoft Metashape to create orthomosaics of the field.}}}
    
% \vspace{0pt}
% \resumeSubheading{MICRO Compiler}{Compilers Course Project}
%     {C, Flex, Bison, Git}{Jan - Apr 2020}
%      {\small{\newline{} {Implemented a compiler for MICRO language. The build of the compiler includes scanning, parsing, symbol table formation, semantic actions, optimizations, and register allocations. It used C-based Flex and Bison utilities}}}
    
% \vspace{0pt}
% \resumeSubheading{Pet Activity Change Detection }{PRML Course Project}
%     {Python, Hidden Markov Models, Deep Learning}{Aug - Nov 2019}
%      {\small{\newline{} {Implemented K-means, Gaussian Hidden Markov Model, and LSTM Neural Networks to detect change in the activities of pet animals based on accelerometer and gyroscopic data. The data came from sensor belts attached to pets and recorded at different instances of time.}}}

% \vspace{0pt}
% \resumeSubheading{Simulation of a Processor}{Computer Architecture Course Project}
%     {Java, OOPS}{Aug - Nov 2019}
%     { {\small{\newline{} {Simulated a pipelined processor in Java. Given an input program in assembly language, the processor will execute the code and produce the results. Analyzed the change in performance of the processor on a set of benchmark programs with different types of caches and pipelines.}}}}

% \vspace{0pt}
% \resumeSubheading{Bulls And Cows Number Guesser}{Data Structures Course Project}
%     {C++, Bash scripting, OOPS}{Jul - Nov 2018}
%     { {\small{\newline{} {Developed a C++ program to guess the "Secret" Number of the opponent team in the game of Bulls \& Cows (A variation of Mastermind). Implemented Min-Max Algorithm and achieved an average of 5.24 moves of guessing the secret number. }}}}
% \resumeSubHeadingListEnd
% \vspace{-5pt}

% \resumeSubHeadingListEnd
% \vspace{-5pt}

\vspace{2pt}


\section{Publications}
\resumeSubHeadingListStart

\resumeSubItem
{Earnings2Insights: Persuasive Investment Report Generation Using Single And Multi-Agent Frameworks (FinNLP @ EMNLP 2025) \href{https://aclanthology.org/2025.finnlp-2.24/}{\faLink}}
{
\newline Co-authored a paper on LLM-based multi-agent frameworks for generating investment reports; ranked 2nd overall and 1st in readability in the Earnings2Insights task at FinNLP.
}

\resumeSubHeadingListEnd

% \section{Skills Summary}
% 	\resumeSubHeadingListStart
% 	\resumeSubItem{Languages}{~~Python, C/C++, JAVA, Matlab, Octave, R, JavaScript, SQL, Bash }
% 	\resumeSubItem{Frameworks}{PyTorch, JAX, Keras, TensorFlow,  HuggingFace, Transformers, Flask, PySpark, Matlpotlib, LlamaFactory}
% 	\resumeSubItem{Tools}{~~~~~~~~~~Docker, GIT, AWS, Oracle DB, Power Bi, Tableau }
%  % resumeSubItem{Languages}{~~~Python, C/C++, JAVA, Matlab, Octave, R, JavaScript, SQL, Bash, PHP, }
% 	% \resumeSubItem{Frameworks}{~Keras, TensorFlow, Flask, Pandas, Scikit, PySpark, Matlpotlib, Plotly, Seaborn, NLTK, PyTorch, RabbitMQ, APIs, Cassandra}
% 	% \resumeSubItem{Tools}{~~~~~~~~~~~Docker, GIT, AWS, Oracle DB, Power Bi, Tableau, Apache Superset, APIs, Google Analytics, Cassandra}
%         \resumeSubHeadingListEnd


\section{Skills Summary}
\resumeSubHeadingListStart

\resumeSubItem{Languages}{Python, C/C++, Java, SQL, R, Bash}

\resumeSubItem{ML / GenAI}{LLMs, Fine-tuning, LLM Evaluations, Synthetic Datasets, Agentic AI, Distributed Training}

\resumeSubItem{Frameworks}{PyTorch, JAX, TensorFlow, HuggingFace Transformers, Keras, PySpark, vLLM}

\resumeSubItem{Systems \& Tools}{Docker, AWS, Deepspeed, Git, Power BI, Tableau}

\resumeSubHeadingListEnd


% % Leadership
% \vspace{-5pt}
% \section{Leadership}
% \resumeSubHeadingListStart
% \resumeSubItem{Secretary, theCodingClub (2019 - 2020) }{Conducted coding workshops, group competitive coding sessions for students at IIT Dharwad. Also organized coding events for Parsec - the institutes first Tech Fest (Feb 2020).}
% % \vspace{2pt}
% \resumeSubItem{Committee member, Rhapsody (2018-2021)}{Founding member of Rhapsody - the Music club of IIT Dharwad. Organized and performed in musical events as the institute's drummer. Conducted Drums lessons for beginner students. }
% % % \vspace{2pt}
% % \resumeSubItem{Member, Career Development Cell, IIT Dharwad (2018-2020)}{Member of CDC - a student-driven body that endeavors to nurture careers of students at the institute. Conducted events and talks to guide the students.}

% \resumeSubHeadingListEnd
% \vspace{-5pt}




% %-----------Position of Responsibilites-----------------
% \section{Position of Responsibilites}
% \resumeSubHeadingListStart
% 	\resumeSubheading
%     {Secretary, theCodingClub}{IIT Dharwad}
%     {IITDh Coding Club}{2019 - 2020}
%     { {\small{\newline{} {Conducted coding workshops for students in the fresher year, had group competitive coding sessions and also hosted Algostrike - a coding competition as an event in IIT Dharwad’s 1st Tech Fest Parsec (Feb 2020) where students all across India participated.}}}}
% \vspace{-3pt}
%     \resumeSubheading
%     {Member, Career Development Cell}{IIT Dharwad}
%     {Support Team, CDC}{2018 - 2020}
%     { {\small{\newline{} {A member of the Support team of CDC - a student-driven body that endeavors to nurture the careers of all the students of the institute. Helped conduct events and talks by professionals in the institute along with helping students excel in their careers.}}}}
% \vspace{-3pt}
%     \resumeSubheading
%     {Committee member, Rhapsody}{IIT Dharwad}
%     {Music Club, IIT Dharwad}{2018 - 2021}
%     { {\small{\newline{} {Was one of the founding members of Rhapsody. As part of the club, conducted and performed in events - musical performances, open mics, concerts, and annual events of national importance. Also conducted Drums learning lessons for beginner students and teachers in the institute. Also participated in music competitions in and outside the campus.}}}}
% \vspace{-3pt}
%     \resumeSubheading
%     {Hostel Secretary}{IIT Dharwad}
%     {Keervani and Abheri Hostel Secretary}{2017 - 2019}
%     { {\small{\newline{} {Was the hostel secretary for Keervani (2017-18) and Abheri (2018-19). Helped resolving student's issue's with their rooms, conducted inter and intra-hostel activites and communicated with the management on a regular basis.}}}}    
% \vspace{-3pt}
%     \resumeSubheading
%     {Student Mentor}{IIT Dharwad}
%     {Student Mentorship Program}{2018 - 2019}
%     { {\small{\newline{} {Was a student mentor and mentored eight 1st year students. Helped resolve their academic queries, guided them to settle in the IITDh campus easily, and provided constructive feedback whenever needed.}}}}    
% \vspace{5pt}



% %-----------Awards-----------------
% \section{Honors and Awards}
% \begin{description}[font=$\bullet$]
% \item {Placed 4th in JioDevelopers Hackathon conducted at the 7th Inter-IIT tech-meet in the Indian Institute of Technology, Bombay }
% \vspace{-5pt}
% \item {Ranked 3024 in the Joint Entrance Examination (JEE) Advanced exam out of over 150k students all across India.}
% \vspace{-5pt}
% \item {Bagged 2nd position in Codex-101, the first coding competition at IIT Dharwad.}


% \end{description}


% %-----------Hobbies-----------------
% \section{Interests and Hobbies}
% \begin{description}[font=$\bullet$]
% \item {I am a drummer and like the Progressive Rock and Heavy Metal genres in Music}
% \vspace{-5pt}
% \item {Really love playing and following Chess.}
% \vspace{-5pt}
% \item {I am an E-sports enthusiast and love playing and following the pro scene in games like Dota 2 and CSGO}
% \vspace{-5pt}
% \item {Love following Japanese media like Anime and Manga.}

% \end{description}

% % \vspace{-5pt}
% % \section{Volunteer Experience}
% %   \resumeSubHeadingListStart
% % 	\resumeSubheading
% %     {Community Lead at Developer Student Clubs NSEC}{Kolkata, India}
% %     {Conducted online and offline technical \& soft-skills training impacting over 3000 students.}{Jan 2019 - Present}
% % \vspace{5pt}
% %     % \vspace{10pt}\textbf{\large{Community Experience}}
% %     \resumeSubheading
% %     {Event Organizer at Google Developers Group Kolkata}{Kolkata, India}
% %     {Organized events, conducted workshops and delivered workshops reaching over 7000 developers.}{Jan 2018 - Present}
% % \resumeSubHeadingListEnd



\end{document}